% chuong_v3/fig_ch09_parabola_peak.tex
% Minh họa hình học phép nội suy đỉnh parabol qua 3 điểm rời rạc (y_-1, y_0, y_+1)
% Trình bày độ lệch fractional lag delta và vị trí cực trị liên tục tau* = tau_p + delta.
\begin{tikzpicture}[
    >=Stealth,
    point/.style={circle, fill=blue!80!black, inner sep=2pt},
    peakpt/.style={circle, fill=red!80!black, inner sep=2.5pt},
    gridline/.style={draw=gray!50, dashed, thick},
    axis/.style={draw=black!80, thick, ->},
    parabola/.style={draw=blue!75!black, very thick}
]

    % Hệ trục tọa độ
    \draw[axis] (-0.5, 0) -- (8.5, 0) node[right, font=\small] {$\tau$ (mẫu trễ)};
    \draw[axis] (0, -0.5) -- (0, 5.2) node[above, font=\small] {$U[\tau]$ (hàm tương quan chuẩn hóa)};

    % Tọa độ các điểm lưới nguyên
    % tau_p - 1 tại x = 2.0, y = 2.8 (y_-1)
    % tau_p     tại x = 4.5, y = 4.5 (y_0)
    % tau_p + 1 tại x = 7.0, y = 3.6 (y_+1) -> y_+1 > y_-1 nên delta > 0 (nghiêng phải)
    \coordinate (P_m1) at (2.0, 2.8);
    \coordinate (P_0)  at (4.5, 4.5);
    \coordinate (P_p1) at (7.0, 3.6);

    % Đỉnh parabol lý thuyết:
    % Với delta = +1/6 trên khoảng cách lưới dx = 2.5 cm:
    % x_peak = 4.5 + (1/6)*2.5 = 4.5 + 0.417 = 4.917 cm
    % y_peak = 4.5 + a*delta^2 + b*delta = 4.5 + 0.07 = 4.57 cm
    \coordinate (P_star) at (4.917, 4.57);

    % Đường dóng thẳng đứng xuống trục hoành
    \draw[gridline] (2.0, 0) -- (P_m1);
    \draw[gridline] (4.5, 0) -- (P_0);
    \draw[gridline] (7.0, 0) -- (P_p1);
    \draw[gridline, draw=red!70!black] (4.917, 0) -- (P_star);

    % Đường dóng nằm ngang sang trục tung
    \draw[gridline] (0, 2.8) -- (P_m1);
    \draw[gridline] (0, 4.5) -- (P_0);
    \draw[gridline] (0, 3.6) -- (P_p1);

    % Vẽ đường cong Parabol đi qua 3 điểm (P_m1, P_0, P_p1)
    % Dạng: y(x) = -0.32*(x - 4.917)^2 + 4.57
    \draw[parabola, domain=1.2:7.8, samples=60] 
        plot (\x, {-0.283*(\x - 4.917)*(\x - 4.917) + 4.57});

    % Điểm rời rạc
    \node[point] at (P_m1) {};
    \node[point] at (P_0)  {};
    \node[point] at (P_p1) {};
    \node[peakpt] at (P_star) {};

    % Nhãn trên trục hoành
    \node[below, font=\scriptsize, fill=white, inner sep=1pt] at (2.0, -0.1) {$\tau_p - 1$};
    \node[below, font=\scriptsize\bfseries, fill=white, inner sep=1pt] at (4.5, -0.1) {$\tau_p$};
    \node[below, font=\scriptsize, fill=white, inner sep=1pt] at (7.0, -0.1) {$\tau_p + 1$};
    \node[below, font=\scriptsize\bfseries, text=red!80!black, fill=white, inner sep=1pt] at (5.05, -0.45) {$\tau^* = \tau_p + \delta$};

    % Nhãn trên trục tung
    \node[left, font=\scriptsize] at (0, 2.8) {$y_{-1}$};
    \node[left, font=\scriptsize] at (0, 4.5) {$y_0$};
    \node[left, font=\scriptsize] at (0, 3.6) {$y_{+1}$};

    % Nhãn tại các điểm
    \node[above left, font=\scriptsize, fill=white, inner sep=1pt] at (P_m1) {$(\tau_p - 1, y_{-1})$};
    \node[above, font=\scriptsize\bfseries, fill=white, inner sep=1.5pt] at (4.3, 4.65) {$(\tau_p, y_0)$};
    \node[above right, font=\scriptsize, fill=white, inner sep=1pt] at (P_p1) {$(\tau_p + 1, y_{+1})$};
    \node[above right, font=\scriptsize\bfseries, text=red!80!black, fill=white, inner sep=1.5pt] at (P_star) {Đỉnh liên tục $\tau^*$};

    % Mũi tên chỉ độ dời delta
    \draw[<->, thick, red!80!black] (4.5, 1.2) -- (4.917, 1.2)
        node[midway, above, font=\scriptsize\bfseries, fill=white, inner sep=1pt] {$\delta$};
    \draw[draw=red!50!black, dashed] (4.5, 0) -- (4.5, 1.3);
    \draw[draw=red!50!black, dashed] (4.917, 0) -- (4.917, 1.3);

    % Khung ghi chú giải thích quy luật dấu
    \node[draw=blue!60!black, fill=blue!3, rounded corners=2pt, font=\scriptsize, align=left, inner sep=4pt] at (3.5, -1.6) {
        \textbf{Quy luật dấu của độ dời đỉnh} $\delta = \frac{y_{+1} - y_{-1}}{2(2y_0 - y_{-1} - y_{+1})}$:\\[2pt]
        $\bullet$ $y_{+1} > y_{-1} \implies \delta > 0$: Đỉnh thực nghiêng sang phải ($\tau^* > \tau_p$),\quad
        $\bullet$ $y_{+1} < y_{-1} \implies \delta < 0$: Đỉnh thực nghiêng sang trái ($\tau^* < \tau_p$),\\[1pt]
        $\bullet$ $y_{+1} = y_{-1} \implies \delta = 0$: Đỉnh thực trùng khít đúng mẫu nguyên $\tau_p$.
    };

\end{tikzpicture}
